\originalchapter{22.02 作业一:核素与量子工具}
题源: Spring 2012 Problem Set 1, 2012年2月16日. 除题中印出的示范外, 本章均为编者独立解答. 核质量取 AME2020 数值, 只列题目所需条目.
\begin{exercise}\label{ex:2202-1-1}
原题 1. 盒内有 147 个中子、95 个质子. (a) 全部组成一个核, 给 N、Z、A、核素符号. (b) 先组成一个 $\alpha$, 余者组成 X, 给两核信息. (c) 先用四质子四中子组成一个核, 余者组成另一个核, 给两核信息.
\end{exercise}
\begin{solution}
(a) $(N,Z,A)=(147,95,242)$, ${}^{242}_{95}$Am. (b) $\alpha={}^4_2$He, 余核 $(145,93,238)$ 即 ${}^{238}_{93}$Np. (c) ${}^8_4$Be 与 $(143,91,234)$ 的 ${}^{234}_{91}$Pa. 注意余核不是 U: 减去四质子后 $95-4=91$, 这也决定下一题的质量条目. ${}^8$Be 不稳定仍是可讨论的核共振, 不因此更改守恒计数.
\end{solution}
\begin{exercise}\label{ex:2202-1-2}
原题 2. (a) 对前题五个核素, 分别用实测质量与半经验公式算总结合能和每核子结合能. (b) 哪种组合在能量上最有利?
\end{exercise}
\begin{solution}
取 $B=[Zm({}^1\mathrm H)+Nm_n-M_{\rm at}]c^2$. 源课的表面系数给一个范围, 为明确数值本章取 $a_v=15.5,a_s=17,a_c=0.691,a_a=23,a_p=34$ MeV, 配对项 $\pm a_pA^{-3/4}$, 使用 $Z(Z-1)$ Coulomb 项. 这是声明的拟合选择, 非声称源题唯一指定系数.
\begin{center}\begin{tabular}{lrrrr}\toprule
核素&AME $B$&AME $B/A$&SEMF $B$&SEMF $B/A$\\\midrule
${}^{242}$Am&1823.468&7.535&1843.076&7.616\\
${}^{238}$Np&1800.761&7.566&1820.230&7.648\\
${}^4$He&28.296&7.074&30.313&7.578\\
${}^{234}$Pa&1777.156&7.595&1796.728&7.678\\
${}^8$Be&56.499&7.062&59.002&7.375\\\bottomrule
\end{tabular}\end{center}
比较三种组合总结合能时, (b) 为 $1829.057$ MeV, (c) 为 $1833.655$ MeV, 三种组合中 (c) 最有利. 正的总结合能增加表示静质量下降; 不含核间 Coulomb 能时该组合更有利, 但还需下一题势垒解释能否快速形成.
\end{solution}
\begin{exercise}\label{ex:2202-1-3}
原题 3. 将前题 (b),(c) 两个核近似成点电荷, 间距 $d=R_1+R_2$, $R_i=1.25A_i^{1/3}$ fm, 计算核间 Coulomb 势能 (MeV).
\end{exercise}
\begin{solution}
$V_C=1.439965Z_1Z_2/[1.25(A_1^{1/3}+A_2^{1/3})]$ MeV. (b) 取 $Z_1Z_2=2\cdot93$, (c) 取 $4\cdot91$. 得约 $27.5$ MeV 与 $51.4$ MeV. 这是接触距离的排斥能, 不是两核已经无穷分离时的 Q 值. 即使分裂能在远处较低, 接触附近仍存在势垒.
\end{solution}
\begin{exercise}\label{ex:2202-1-4}
原题 4. 八个操作为 $x\phi$, $-\ii\hbar\phi'$, $-\hbar^2\phi''/(2m)$, $I\phi=\phi$, $F\phi=F(x)\phi$, $A\phi=0$, $B\phi=\phi/21$, $C\phi=33$. (a) 求各自平方与逆 (若存在). (b) 哪些是线性的?
\end{exercise}
\begin{solution}
平方依次为 $x^2\phi$, $-\hbar^2\phi''$, $\hbar^4\phi^{(4)}/(4m^2)$, $\phi$, $F^2\phi$, $0$, $\phi/441$, $33$. $I^{-1}=I$, $B^{-1}=21I$; $x,F$ 的逆分别乘 $1/x,1/F$, 要求所选函数空间及零点允许. 零算符和常数映射都不单射, 无逆. 导数算符有常数核, 二阶导数有仿射函数核, 因而在一般光滑函数空间没有双边逆; 给边界条件并排除零模后才能用 Green 函数求逆. 前七个都满足线性性, 包括零算符; C 不满足 $C0=0$, 非线性. 原题名称``annihilation operator''这里指零映射, 不是振子降低算符.
\end{solution}
\begin{exercise}\label{ex:2202-1-5}
原题 5 (含已解示范). 由 $p=\hbar k=mv$ 推出非相对论 de Broglie 波长 $\lambda=2\pi\hbar/\sqrt{2mE}$. (a) 晶格间距几 Å, 选 1 MeV 快中子还是 0.025 eV 热中子用于散射? 数值计算两波长. (b) C$_{60}$ 速度 220 m/s, 波长是多少? 同 50 nm 宽、100 nm 周期的光栅比较.
\end{exercise}
\begin{solution}
源示范重编: $k=2\pi/\lambda$, $p^2=2mE$, 两式消去 p 即得到公式. 用 $m_nc^2\simeq939.565$ MeV, 快中子 $\lambda\simeq28.6$ fm, 热中子 $\lambda\simeq1.81$ Å, 所以后者符合晶格尺度. C$_{60}$ 质量约 $720\,\mathrm u$, $\lambda=h/(Mv)\simeq2.52$ pm. 它远小于光栅周期但衍射角约 $\lambda/d\simeq2.5\times10^{-5}$ rad, 可通过足够传播距离分辨, 不要求波长等于缝宽才能发生衍射.
\end{solution}
\begin{exercise}\label{ex:2202-1-6}
原题 6. (a) 定义本征值方程并给两个例子. (b) 求平面旋转 $R(\theta)=\bigl(\begin{smallmatrix}\cos\theta&-\sin\theta\\\sin\theta&\cos\theta\end{smallmatrix}\bigr)$ 在 $\theta=\pi/3$ 时本征值与向量. (c) 与 x 轴夹 $60^\circ$ 的向量分别在原坐标、转过 $45^\circ$ 的基中表示. (d) 对力 $F=-kx$ 写一维经典振子方程并解释、求相应本征值与本征函数.
\end{exercise}
\begin{solution}
(a) $O\phi=\lambda\phi$, $\phi\ne0$; 例如 $H\phi=E\phi$, $-\ii\partial_x\e^{\ii kx}=k\e^{\ii kx}$. (b) 在复空间本征值 $\e^{\pm\ii\pi/3}$, 对应 $(1,\mp\ii)^T/\sqrt2$; 实空间无非零本征向量. (c) 向量长度 r 时原分量 $r(1/2,\sqrt3/2)$, 旋转基中为 $r(\cos15^\circ,\sin15^\circ)$. (d) $\ddot x=-(k/m)x$, 可把 $-\dd^2/\dd t^2$ 的本征值视为 $\omega^2=k/m$, 本征函数 $\e^{\pm\ii\omega t}$ 或 sin、cos. 给定初始条件才确定具体振动; 未给时间边界时不是离散的谱问题.
\end{solution}
\begin{exercise}\label{ex:2202-1-7}
原题 7. 证明 (a) $\delta(x)=\delta(-x)$ (原题已解), (b) $\delta(ax)=\delta(x)/a$, (c) $f(x)\delta(x-x_0)=f(x_0)\delta(x-x_0)$.
\end{exercise}
\begin{solution}
(a) 对测试函数积分作 $u=-x$ 得同一 $f(0)$. (b) 原题漏绝对值: $a\ne0$ 时正确式为 $\delta(ax)=\delta(x)/|a|$; 换元后积分测度给绝对雅可比, $a<0$ 不能得到负 delta. $a=0$ 没有这种合法分布缩放. (c) 两侧作用于任意测试函数 g 都等于 $f(x_0)g(x_0)$, 要求 f 在该点足够连续.
\end{solution}
\begin{exercise}\label{ex:2202-1-8}
原题 8. 用 $\delta(x)=(2\pi)^{-1}\int\e^{\ii kx}\dd k$ 求 $\cos(k_0x)$ 的 Fourier 变换.
\end{exercise}
\begin{solution}
采用第 \ref{ch:qm-core} 章对称约定, 将 cos 拆为两指数, 得 $\widetilde f(k)=\sqrt{\pi/2}[\delta(k-k_0)+\delta(k+k_0)]$. 两峰代表正反传播动量, 不表示两个不同动能.
\end{solution}
